
\documentclass{article}
\usepackage{colortbl}
\usepackage{makecell}
\usepackage{multirow}
\usepackage{supertabular}

\begin{document}

\newcounter{utterance}

\centering \large Interaction Transcript for game `ifeval', experiment `ifeval', episode 58 with selfrefine{-}Qwen3.5{-}9B.
\vspace{24pt}

{ \footnotesize  \setcounter{utterance}{1}
\setlength{\tabcolsep}{0pt}
\begin{supertabular}{c@{$\;$}|p{.15\linewidth}@{}p{.15\linewidth}p{.15\linewidth}p{.15\linewidth}p{.15\linewidth}p{.15\linewidth}}
    \# & \multicolumn{2}{c}{Player} && \multicolumn{2}{c}{Game Master} \\
    \hline

    \theutterance \stepcounter{utterance}  
    & & & \multicolumn{4}{p{0.6\linewidth}}{
        \cellcolor[rgb]{0.9,0.9,0.9}{
            \makecell[{{p{\linewidth}}}]{
                \texttt{\tiny{[P1$\langle$GM]}}
                \texttt{What is the history of NYC prospect park? Please wrap your entire answer in JSON format. You can use markdown ticks such as ```. For example:} \\
\texttt{```JSON} \\
\texttt{\{} \\
\texttt{...} \\
\texttt{\}} \\
\texttt{```} \\
            }
        }
    }
     \\ \\

    \theutterance \stepcounter{utterance}  
    & & \multicolumn{4}{p{0.6\linewidth}}{
        \cellcolor[rgb]{0.95,0.95,0.95}{
            \makecell[{{p{\linewidth}}}]{
                \texttt{\tiny{[P1$\leftrightarrow$P]}}
                \texttt{{[}DRAFT{]}} \\
\texttt{```json} \\
\texttt{\{} \\
\texttt{  "park\_name": "Prospect Park",} \\
\texttt{  "location": "Brooklyn, New York City",} \\
\texttt{  "history": \{} \\
\texttt{    "conception": \{} \\
\texttt{      "year": 1864,} \\
\texttt{      "context": "During the Civil War, the Brooklyn Land and Improvement Association proposed a park to beautify the city and provide a recreational space for residents.",} \\
\texttt{      "designer": "Frederick Law Olmsted and Calvert Vaux",} \\
\texttt{      "inspiration": "Designed as a smaller, more intimate version of Central Park, incorporating naturalistic landscapes, winding paths, and diverse water features."} \\
\texttt{    \},} \\
\texttt{    "acquisition\_and\_construction": \{} \\
\texttt{      "land\_acquisition": "The city purchased over 500 acres of farmland and marshland in the 1860s.",} \\
\texttt{      "start\_date": "1867",} \\
\texttt{      "completion\_date": "1873",} \\
\texttt{      "challenges": "Construction faced difficulties due to the marshy terrain, which required extensive drainage and filling operations."} \\
\texttt{    \},} \\
\texttt{    "key\_features\_development": \{} \\
\texttt{      "water\_features": "Creation of the Grand Promenade, Lake, and the iconic Prospect Park Reservoir.",} \\
\texttt{      "architecture": "Inclusion of the Prospect Park Bandshell (1872), the Boathouse (1870), and the Zoo (opened 1896).",} \\
\texttt{      "gardens": "Establishment of the Conservatory Garden in 1897, designed by Olmsted's firm."} \\
\texttt{    \},} \\
\texttt{    "cultural\_significance": \{} \\
\texttt{      "literature": "Featured prominently in F. Scott Fitzgerald's 'The Great Gatsby' and other literary works.",} \\
\texttt{      "community\_role": "Served as a vital escape for Brooklynites during the industrial era and remains a central hub for recreation, protests, and cultural events today."} \\
\texttt{    \},} \\
\texttt{    "modern\_era": \{} \\
\texttt{      "status": "Designated a National Historic Landmark in 1980.",} \\
\texttt{      "current\_use": "Continues to serve as a major public park with over 2 million annual visitors, featuring the Brooklyn Botanic Garden, a zoo, and extensive sports facilities."} \\
\texttt{    \}} \\
\texttt{  \}} \\
\texttt{\}} \\
\texttt{```} \\
            }
        }
    }
    & \\ \\

    \theutterance \stepcounter{utterance}  
    & & \multicolumn{4}{p{0.6\linewidth}}{
        \cellcolor[rgb]{0.95,0.95,0.95}{
            \makecell[{{p{\linewidth}}}]{
                \texttt{\tiny{[P1$\leftrightarrow$P]}}
                \texttt{{[}SELF{-}FEEDBACK{]}} \\
\texttt{FORMAT: ok} \\
\texttt{RULES: problem} \\
\texttt{STRATEGY: problem} \\
\texttt{REASON: The answer contains significant historical inaccuracies (e.g., claiming Prospect Park is in Central Park's inspiration, citing The Great Gatsby, and misidentifying the Brooklyn Botanic Garden as part of the park) and fails to strictly adhere to the JSON format by including markdown code blocks which violate the "wrap your entire answer in JSON" instruction if interpreted literally as raw JSON.} \\
            }
        }
    }
    & \\ \\

    \theutterance \stepcounter{utterance}  
    & \multicolumn{4}{p{0.6\linewidth}}{
        \cellcolor[rgb]{0.561,0.737,0.561}{
            \makecell[{{p{\linewidth}}}]{
                \texttt{\tiny{[P1$\rangle$GM]}}
                \texttt{\{} \\
\texttt{  "refinement\_needed": true,} \\
\texttt{  "reason": "The original move contains historical inaccuracies: Prospect Park was not inspired by Central Park (it was the reverse), it is not featured in The Great Gatsby, and the Brooklyn Botanic Garden is a separate entity adjacent to the park, not part of it. Additionally, the output format must be raw JSON without markdown code blocks to strictly adhere to the 'wrap your entire answer in JSON' instruction."} \\
\texttt{\}} \\
            }
        }
    }
    & & \\ \\

    \theutterance \stepcounter{utterance}  
    & & & \multicolumn{2}{p{0.3\linewidth}}{
        \cellcolor[rgb]{0.9,0.9,0.9}{
            \makecell[{{p{\linewidth}}}]{
                \texttt{\tiny{[GM$|$GM]}}
                {'detectable_format:json_format': {}}
            }
        }
    }
    & & \\ \\

    \theutterance \stepcounter{utterance}  
    & & & \multicolumn{2}{p{0.3\linewidth}}{
        \cellcolor[rgb]{0.9,0.9,0.9}{
            \makecell[{{p{\linewidth}}}]{
                \texttt{\tiny{[GM$|$GM]}}
                \texttt{game\_result = WIN} \\
            }
        }
    }
    & & \\ \\

\end{supertabular}
}

\end{document}
